\section{Training, reference generation, and active learning}

\subsection{Hierarchical teacher stack}

Training uses a hierarchy rather than one monolithic truth solver:
\begin{description}
\item[Level 0:] analytic filament, mutual-inductance, homogeneous-medium, and
  simple thermal Green estimates;
\item[Level 1:] low-order integral-equation/PEEC solve;
\item[Level 2:] converged hybrid conductor/SIE/VIE electromagnetic teacher and
  thermal boundary/Green teacher;
\item[Level 3:] independent external or legacy high-fidelity solver used on a
  small audit subset for cross-method validation.
\end{description}

A datum records the continuous scene and teacher policy, not a mesh-dependent
feature vector.

\subsection{Dataset unit}

One immutable scene sample contains
\[
 \mathcal D_j=
 \{\mathcal S_j,\Omega_{\omega,j},\Omega_{T,j},
 Z_j,\mathcal R_{e,j},\mathcal Q_j,\mathcal R_{T,j},
 \text{certificates}\}.
\]
Here $\mathcal R_e$ is a structured electrical realization, $\mathcal Q$ the
continuous/low-rank loss representation, and $\mathcal R_T$ the thermal reduced
realization. Quadrature data may be cached internally but is not part of the
model schema.

\subsection{Split policy}

Validation and test scenes are frozen before active learning. Splits are grouped
by geometric/material regimes, not by random neighboring samples. Additional
active-learning scenes enter the training pool only. A completely independent
release suite remains untouched by model selection.

\subsection{Losses}

Because major physical properties are hard-parameterized, the trainable loss
focuses on observable accuracy:
\[
 \mathcal L
 =
 \lambda_Z\mathcal L_Z
 +\lambda_{\mathrm{dyn}}\mathcal L_{\mathrm{dyn}}
 +\lambda_q\mathcal L_q
 +\lambda_T\mathcal L_T
 +\lambda_r\mathcal L_{\mathrm{res}}
 +\lambda_u\mathcal L_{\mathrm{uncert}}.
\]
Matrix errors use physically scaled norms and current-space probes. Spatial loss
errors are measured after decoding and power normalization. Dynamic losses
compare transfer functions, impulse responses, and selected time trajectories.

\subsection{Physics-residual active learning}

Candidate scenes are not selected primarily by Euclidean parameter distance.
For a candidate $\mathcal S$, FAST mode provides a prediction and cheap
estimators:
\[
 \eta(\mathcal S)
 =
 w_e\eta_{\mathrm{EM}}
 +w_p\eta_{\mathrm{power}}
 +w_t\eta_{\mathrm{thermal}}
 +w_u\eta_{\mathrm{uncertainty}}.
\]
High-score scenes are promoted to higher teacher levels. In CERTIFIED mode the
actual integral residual
\[
 \eta_{\mathrm{EM}}
 =
 \frac{\norm{\mathcal B i-\mathcal A x_\theta}}
      {\norm{\mathcal B i}}
\]
is directly available for representative port excitations.

\subsection{Curriculum}

A recommended curriculum is:
\begin{enumerate}
\item one and two coils in homogeneous background;
\item arbitrary relative pose and continuous superellipse geometry;
\item finite-cross-section skin/proximity effects;
\item passive packages and piecewise homogeneous media;
\item localized inhomogeneous/anisotropic media;
\item frequency and temperature dependence;
\item full electrothermal trajectories and external circuits.
\end{enumerate}
Each stage retains previous-stage release tests to prevent catastrophic
regression.

\subsection{No fixed truth count}

Numbers such as 96, 256, or 1024 are not theoretical requirements. Training
terminates when fixed validation regimes and independent residual certificates
meet the declared error budget or when a declared resource ceiling is reached.
